\section{Related works}
\label{sec:related_works}
\noindent\textbf{Human and object reconstruction.}
Most recent works on human and object reconstruction~\cite{chen2019holistic++,zhang2020perceiving,xie2022chore,nam2024contho,xie2023templatefree,dwivedi2025interactvlm,wen2025reconstructing,li2025scorehoi,xu2021d3d,xie2023visibility,xue2026rhino,xie2025intertrack,xie2026cari4d,wen2025efficient,huang2026recovering} reconstruct the 3D spatial relationship between humans and objects from a single image or monocular video, relying on learned interaction priors, such as contact maps.
CONTHO~\cite{nam2024contho} jointly predicts dense contact maps from images and refines the reconstruction using a contact-based Transformer.
VisTracker~\cite{xie2023visibility} exploits the correlation between human and object motion, predicting the object pose of occluded frames from the human motion and the neighboring visible frames.
RHINO~\cite{xue2026rhino} jointly recovers the human, object, and scene from a monocular video, and refines their poses with contact-driven optimization.
CARI4D~\cite{xie2026cari4d} employs a render-and-compare network that compares rendered estimates against input observations to enhance reconstructions.


Unlike existing 4D human and object reconstruction methods, our framework has two key differences.
First, while prior works rely primarily on local geometric priors (\textit{e.g.}, contact maps) without high-level semantic understanding, we incorporate semantic text guidance to capture complex real-world interactions. By temporally grounding these descriptions into distinct sub-actions, we provide precise, phase-aware contexts that local geometry alone cannot represent.
Second, whereas previous methods reconstruct solely from monocular observations and leave severe mutual occlusion unresolved, we leverage our semantic prompts to synthesize multi-view videos. This explicitly reveals hidden regions that single-view approaches miss, providing reliable and physically consistent auxiliary evidence for the 4D optimization.


\noindent\textbf{Text guidance for reconstruction.}
Recently, various 3D reconstruction methods~\cite{deng2023nerdi,huang2024tech,tang2023make,ye2024ghop,tang2024dreamgaussian,sun2024dreamcraft3d,qian2024magic123,wang2024animatabledreamer,nam2026tehor,kim2025david} have leveraged textual information as semantic guidance through text-to-image diffusion models.
NeRDi~\cite{deng2023nerdi} reconstructs a 3D object from a single image by using text captioning to align novel-view appearances with the caption via a text-to-image diffusion model~\cite{rombach2022high}.
TeCH~\cite{huang2024tech} extracts detailed text prompts from the input image about the garments, colors, and hairstyles of a clothed human, and uses them to reconstruct the 3D clothed human.
TeHOR~\cite{nam2026tehor} leverages text descriptions to enhance joint reconstruction of 3D human and object by capturing semantic interaction cues.
Unlike these methods, which use a single text description for static 3D reconstruction, our framework extends text guidance to 4D reconstruction by capturing the temporal variation of human-object interaction.
To this end, we divide the input video into semantically distinct sub-actions and assign a specific text description to each segment, providing temporally aligned semantic guidance for 4D reconstruction.


\noindent\textbf{Video temporal grounding.}
Video temporal grounding methods localize the segment of an input video that semantically corresponds to a text query~\cite{gao2017tall,zhang2020learning,lei2021moment,moon2023query,lin2023univtg,ren2024timechat,huang2024vtimellm,wang2024groundedvideollm,chen2025fmimic,zhang2026timelens}.
TALL~\cite{gao2017tall} introduces a cross-modal temporal regression method to predict temporal action boundaries aligned with a text query.
UniVTG~\cite{lin2023univtg} proposes a unified temporal grounding framework that handles diverse temporal label formats, such as timestamps and video intervals.
TimeLens~\cite{zhang2026timelens} adapts pretrained multimodal large language models~\cite{bai2025qwen3} with interleaved timestamp encoding and reinforcement learning for precise temporal grounding.
FMimic~\cite{chen2025fmimic} performs temporal grounding based on contact distance thresholds to identify human--object interaction states for robotic imitation learning.
Different from FMimic, our framework uses a coarse-to-fine strategy that detects initial sub-action boundaries from object trajectories and refines them into precise temporal spans using TimeLens~\cite{zhang2026timelens}.
This enables semantically meaningful temporal chunking aligned with detailed text queries, establishing accurate phase-specific boundaries for 4D HOI reconstruction.